% !PRINTED 662
\weberpage{662}
It is to be proved that this degree \gesperrt{cannot be less than $n_1$}, or, what
means the same thing, that $\eta_1, \eta_2, \ldots, \eta_{n_1}$ admit of
a rational representation by $\tau_1$. This follows if one
differentiates (10) partially with respect to one of the $t_\nu$:
\[
\Phi'(\tau_1)\,\eta_\nu + \frac{\partial A_1}{\partial t_\nu}\tau_1^{m-1}
+ \cdots + \frac{\partial A_{m-1}}{\partial t_\nu}\tau_1
+ \frac{\partial A_m}{\partial t_\nu} = 0.
\]

Now since $\Phi'(\tau_1)$ does not vanish when $m$ is the lowest
possible degree of the equation (10), one obtains from this at once the
required representation of $\eta_\nu$. Hence $m = n_1$ must hold. With
this it is proved:

2. \gesperrt{The functional $\tau_1$ satisfies an irreducible equation
of degree $n_1$ in $z_1$.}

If therefore we put $r = 0, 1, 2, \ldots, n_1-1$ in (9), there results a
system of linear equations for $\eta_1, \eta_2, \ldots, \eta_{n_1}$
whose determinant is a non-vanishing functional in $z_1$, and one can
therefore assign to the variable $t_i$ such constant values that the
determinant does not vanish even then. If $\tau_1$ passes by this
substitution into $\theta_1$, then from the system (9) one can express
$\eta_1, \eta_2, \ldots, \eta_{n_1}$ rationally by $z_1$ and $\theta_1$.

With this, however, it is proved:

3. \gesperrt{If $z_1$ is an arbitrary, non-constant function in
$\Omega$, and $\theta_1$ a second function of the same field which
satisfies no equation of lower degree than $n_1$ in $z_1$, then all
functions of the field $\Omega$, and hence also $z$ and $\theta$, are
rationally expressible by $z_1$ and $\theta_1$.}

Between any two functions $\alpha, \beta$ of the field $\Omega$ there
always subsists an equation with constant coefficients

\weq{11}{F(\alpha, \beta) = 0,}
and among all possible equations of this form there is one of the lowest
possible degree, both with respect to $\alpha$ and with respect to
$\beta$.

\begin{center}
\rule{4cm}{0.4pt}
\end{center}
